\subsection{Physical system and reference frame}

We describe the motion in a 2-D frame (Fig.~\ref{fig:geometry}).
The $x$ axis points away from Earth, $y$ points along the direction of
flight, and $z$ (out of the plane) is perpendicular to the orbit, representing the spin axis.

The satellite (body and panels) is a rigid body with mass $M (kg)$ and moment
of inertia $I (kg\cdot m^2)$ about the $z$ axis through its centre of mass (CoM). In a frame
moving with the satellite, it is at rest before the hit.

The fragment is a point mass $m (kg)$ moving at speed $v_0 (m/s)$. It hits
the panel at the point $\vect r (m)$ (measured from the CoM), at an angle
$\theta$ from the outward normal $\hat{\vect n}$ of the panel. With the
panel tangent $\hat{\vect t}=\hat{\vect z}\times\hat{\vect n}$, the
velocity of the fragment and its speed along the normal are:
\begin{equation}
  \vect v_{d0}=-v_0\left(\cos\theta\,\hat{\vect n}+\sin\theta\,\hat{\vect t}\right),
  \qquad v_n=v_0\cos\theta .
  \label{eq:vd0}
\end{equation}
So $\theta=0$ is a head-on hit and $\theta$ close to $90^\circ$ is a
grazing hit.

\begin{itemize}
  \item $\vect v_{d0}$: the initial velocity  of the fragment $(m/s)$
  \item $v_n$: the speed of the fragment along the normal $(m/s)$
\end{itemize}

\begin{figure}[H]
  \centering
  \begin{tikzpicture}[>=Stealth, font=\small]
    \draw[thick, fill=gray!20] (-0.7,-0.55) rectangle (0.7,0.55);
    \draw[thick, fill=blue!15] (0.7,-0.08) rectangle (6.2,0.08);
    \fill (0,0) circle (1.6pt);
    \node[above] at (0,0.02) {CoM};
    \node[below] at (0,-0.55) {satellite $M,\ I$};
    \draw[->, thick] (0,0) -- (5.2,0);
    \node[below] at (3.0,-0.08) {$\vect r$};
    \coordinate (P) at (5.2,0.08);
    \draw[->, thick, red!70!black] (P) -- (5.2,1.5) node[above] {$\hat{\vect n}$};
    \draw[->, thick, red!70!black] (P) -- (4.3,0.08);
    \node[above, red!70!black] at (4.4,0.08) {$\hat{\vect t}$};
    \draw[->, very thick, orange!80!black] (5.95,1.38) -- (5.27,0.2);
    \fill[orange!80!black] (5.95,1.38) circle (2.2pt);
    \node[right, orange!80!black] at (6.0,1.38) {debris $m$, $\vect v_{d0}$};
    \draw (5.2,0.78) arc[start angle=90, end angle=60, radius=0.7];
    \node at (5.42,1.05) {$\theta$};
    \draw[->] (-3.2,-0.5) -- (-2.5,-0.5) node[right] {$x$};
    \draw[->] (-3.2,-0.5) -- (-3.2,0.2) node[above] {$y$};
    \node[below left] at (-3.2,-0.5) {$z\,\odot$};
  \end{tikzpicture}
  \caption{Geometry of the impact: a fragment of mass $m$ hits the panel at
  lever arm $\vect r$ from the satellite's centre of mass, at angle
  $\theta$ from the panel normal $\hat{\vect n}$.}
  \label{fig:geometry}
\end{figure}
When the fragment hits the panel, it transfers momentum to the satellite. This can change both the satellite's linear and rotational motion. If the impact happens away from the centre of mass, it can also create a torque, causing the satellite to rotate. The impact angle affects how much of the fragment's velocity acts in the normal and tangential directions.
